\usepackage{fancyhdr}

% Dimensiones de la página.
\geometry{paper=a4paper, top=2cm, left=2cm, right=2cm}


\pagestyle{fancy}

% Esconde la numeración por defecto de las páginas.
\fancyhf{}

% Logos
\lhead{\includegraphics[width=2.5cm]{imagenes/logo_usm}}
\rhead{\includegraphics[width=2.5cm]{imagenes/logo_ciac}}

%\chead{
%  Evaluación formativa \textit{Online}  N° 3\\
%  IWI-131 Semestre II-2019 Fase II \\
%  CIAC Casa Central
%}
% El titulo de la guía/pauta
\newcommand{\tituloDocumento}[1]{\chead{#1}}
\newcommand{\tituloGuia}[2]{\chead{
  #1 Guía N° #2 \\
  \ramo \ Semestre \ \semestre-\anno \\
  CIAC Casa Central
}}
\newcommand{\tituloPauta}[2]{\chead{
  #1 Pauta N° #2 \\
  \ramo \ Semestre \ \semestre-\anno \\
  CIAC Casa Central
}}
\newcommand{\tituloOtro}[1]{\chead{
  #1 \\
  \ramo \ Semestre \ \semestre-\anno \\
  CIAC Casa Central
}}

\newcommand{\ramo}{\textbackslash renewcommand\{\textbackslash ramo\}\{IWI-131\}}

\newcommand{\anno}{\textbackslash renewcommand\{\textbackslash anno\}\{2020\}}
\newcommand{\semestre}{\textbackslash renewcommand\{\textbackslash semestre\}\{I\}}


% Debes definir el comando \anno (\newcommand{\anno}{2020})
\lfoot{Página \thepage}
\rfoot{\LaTeXe / CIAC \anno}
% \rfoot{\LaTeXe / CIAC 2020}

% Tamaño de las lineas horizontales al inicio y final de la página
\renewcommand{\headrulewidth}{0.5pt}
\renewcommand{\footrulewidth}{0.5pt}

% Deja espacio suficiente para que el contenido no choque con el encabezado
\setlength{\headheight}{47pt}

% Cambia la fuente del comando \texttt
\renewcommand{\ttdefault}{pcr}


